\section{Limitations}
\label{sec:limitations}

\paragraph{Controller expansion, fixed but not yet ahead.} All five of MAGE-RAG's relation types are instantiated (\S\ref{sec:graph}); the controller's frontier expansion (Algorithm~\ref{alg:controller}) originally followed only \code{reading\_order}, which never crosses a page boundary, and ablation (c) (\S\ref{sec:experiments}) measured the resulting gap against static top-$k$ on cross-page questions. Adding \code{links\_to}/\code{continues} expansion closes most of it (29.7\%~$\to$~39.6\% at default seed width) but the controller still does not strictly beat the baseline: \code{layout\_adjacency} and \code{section\_hierarchy} remain unconsumed, and \code{semantic\_neighbor} is off by default.

\paragraph{Single-hop self-generated data cannot exercise the graph -- addressed via HotpotQA.} \S\ref{sec:contrastive}'s QA generation asks for a question answerable only from the tile shown, so an oracle controller could answer every such question from one node. \code{src/multihop\_qa.py} closes this gap with HotpotQA~\citep{hotpotqa2018} instead: each question's \code{supporting\_facts} names two Wikipedia articles and the sentences in each needed to answer it, so we render both, keep tiles overlapping a supporting sentence as positives, and mine hard negatives by TF-IDF pooled across both articles' tiles. Scaled to 91/150 sampled questions (61\%) and used in ablation (c). Using this data for training, rather than evaluation only, remains future work.

\paragraph{Chrome removal is Wikipedia-specific; tiling itself is not.} All quantitative results here are on Wikipedia, and chrome removal (\S\ref{sec:tiling}) uses a MediaWiki-specific selector list. We rendered two non-Wikipedia pages (Python's documentation, MDN) as a qualitative check: the tiling mechanism itself (bounding boxes, merging, patching) produces coherent tiles over each page's article body, but each site's own navigation/sidebar chrome survives untiled, since it matches no selector we remove -- what is Wikipedia-specific is the chrome list, not the tiling logic, though we make no quantitative claim beyond this observation. JS-hydrated SPAs and \code{div}-soup layouts remain out of scope, the regime web-agent work such as WebArena~\citep{webarena2023}, Mind2Web~\citep{mind2web2023}, and SeeAct~\citep{seeact2024} already studies from the acting side.

\paragraph{Other gaps.} No result above is at the level of the final \emph{answer}: ablations (a)/(b) score tile retrieval (ranking, R@1/MRR) and ablation (c) scores evidence coverage (whether the right tiles were opened), never whether the pipeline's synthesized answer is actually correct end-to-end. \code{scripts/eval\_multihop\_e2e.py} closes this gap (evidence controller output fed to the reader for answer synthesis, scored by EM/F1 against HotpotQA gold answers, same normalization as SQuAD/HotpotQA). Run on $n{=}20$ of the 91 multi-hop questions of ablation (c) -- a random subset (\code{random.Random(seed{=}0).sample}, disclosed and reproducible, not a dataset-order prefix), kept small given the $\sim$2\,min/tile VLM cost per opened tile plus one further call for answer synthesis -- it gives EM~$=$~0.050, F1~$=$~0.057: far below ablation (c)'s $\sim$40\% evidence-coverage rate, so even when the controller opens a correct tile, final answer synthesis from that evidence still fails most of the time on this sample. This end-to-end gap is not yet diagnosed (mis-read evidence text vs.\ an under-specified synthesis prompt vs.\ the base VLM's own weak reading ability, cf.\ ablation (b)) and the small sample means this estimate itself carries wide uncertainty; running the remaining 71 questions once triaged is the natural next step. Relatedly, ablation (b)'s gain has no external anchor yet: synthetic questions, the optional false-negative judge, and the fine-tuned reader all share the \texttt{qwen3-vl}/\texttt{Qwen2-VL} family, so the held-out Recall@1 improvement is entirely on our own model-generated questions; \code{scripts/eval\_hotpotqa\_reader.py} reruns it against this section's human-authored HotpotQA questions instead, implemented but not yet run to completion. Patching (\S\ref{sec:tiling}) splits an oversized element into equal-height sub-tiles regardless of content, a narrower version of PixelRAG's own failure mode. The controller's TF-IDF relevance scoring (\S\ref{sec:controller}) misses paraphrases and cross-modal cues the fine-tuned reader would catch; wiring in the reader's own similarity is the natural next iteration. Hard-negative mining is single-page outside HotpotQA; the corpus crawl is one hop, non-recursive; hyperparameters (Appendix~\ref{sec:appendix-hyperparams}) are chosen once, not swept.
